% Path Complexity Metrics - LaTeX Table for Publication
% Auto-generated by Step 6

\begin{table}[htbp]
\centering
\caption{\textbf{Path complexity metrics for cohort trajectories (2007--2024).} 
Tortuosity measures the ratio of actual path length to straight-line displacement; 
values near 1.0 indicate direct movement while higher values reflect meandering or 
period-specific fluctuations. The mean resultant length (R) quantifies directional 
consistency, with values near 1.0 indicating all period-to-period transitions point 
in similar directions. Circular standard deviation provides an intuitive measure of 
angular spread.}
\label{tab:path_complexity}

\begin{tabular}{lccccc}
\toprule
\textbf{Age Group} & 
\textbf{Displacement} & 
\textbf{Path Length} & 
\textbf{Tortuosity} & 
\textbf{R} & 
\textbf{Circular SD} \\
 & (km) & (km) & (ratio) & & (degrees) \\
\midrule
10-17 & 0.539 & 2.222 & 4.13 & 0.203 & 102.3 \\
18-30 & 0.510 & 2.816 & 5.53 & 0.472 & 70.2 \\
31-55 & 0.448 & 2.094 & 4.67 & 0.195 & 103.5 \\
56-65 & 0.064 & 3.466 & 54.58 & 0.587 & 59.2 \\
66+ & 0.146 & 2.117 & 14.48 & 0.192 & 104.2 \\
\bottomrule
\end{tabular}

\vspace{0.3cm}
\begin{tablenotes}
\small
\item \textit{Displacement}: Straight-line distance from baseline (2007--09) to final period (2022--24).
\item \textit{Path Length}: Total distance traveled through action-space across all period-to-period transitions.
\item \textit{Tortuosity}: Path length / displacement. Values of 1.0 indicate perfectly straight trajectories; 
higher values indicate deviation from direct movement.
\item \textit{R (Mean Resultant Length)}: Circular statistic measuring directional consistency (range: 0--1). 
Higher values indicate period-to-period vectors point in more similar directions.
\item \textit{Circular SD}: Angular dispersion of trajectory directions in degrees. 
Lower values indicate tighter angular clustering.
\end{tablenotes}
\end{table}

